\chapter{Robot Platform and Propulsors}
\label{ch:platform}

This chapter describes the quadruped BODY2, the two propulsors that are compared and the leg motions that
drive them. The motions defined here are the inputs of the CFD in \cref{ch:cfd,ch:results}.

\section{The Quadruped BODY2}
\label{sec:plat-robot}

BODY2 is a small amphibious quadruped built at the chair. Its sealed, flat-bottomed hull carries the
electronics and the eight servos and gives the robot its buoyancy (\cref{fig:robot}). At the design waterline
the hull is \SI{232}{\milli\metre} long and \SI{99}{\milli\metre} wide, with a draft of \SI{25}{\milli\metre}
and a displacement of about \SI{385}{\gram}. The complete robot weighs \SI{418.8}{\gram} and therefore floats \SIrange{1}{2}{\milli\metre}
deeper, consistent with the waterline observed on the hardware. The robot swims at the surface: the hull floats and the legs
reach down into the water.

Each leg is a planar linkage with two degrees of freedom that moves in a vertical plane parallel to the
hull's symmetry plane. Two servos (PTK\,7465W; stall torque \SIrange{0.44}{0.57}{\newton\metre} and no-load speed \SIrange{650}{860}{\degree\per\second} at \SIrange{6.0}{8.4}{\volt}) drive it. The first rotates a crank about the hip axis $O_2$ (angle $\psi$);
the second sets the configuration of a parallelogram that carries the foot (angle $\phi$). The end point $E$
of the parallelogram carries the propulsor. For walking, the foot is the fan paddle described below. For
swimming, the leg can be fitted either with this paddle or with a fluke on the same mounting points, so that
the two propulsion principles can be compared on one leg.

\begin{figure}[t]
  \centering
  \includegraphics[width=\textwidth]{fig_robot}
  \caption[BODY2 with fan paddles and with flukes]{BODY2 in its two swimming configurations, rendered from the
    CAD models. (a) Drag-type configuration with the folding fan paddles (green) used for walking and
    paddling. (b) Lift-type configuration with a fluke (blue) on a straight foot rod at each leg.}
  \label{fig:robot}
\end{figure}

\section{Propulsors}
\label{sec:plat-propulsors}

\Cref{fig:propulsors} and \cref{tab:propulsors} summarize the two propulsors.

\paragraph{Drag-type: folding fan paddle.} The foot is a \SI{120}{\degree} fan of radius
\SI{42}{\milli\metre} at the end of a \SI{50}{\milli\metre} rod; \emph{fan} denotes the folding blade alone and \emph{paddle} the
blade on its rod. Open, the fan has an area of
\SI{18.5}{\square\centi\metre} (\SI{20.5}{\square\centi\metre} including the rod). A cable folds it for the
recovery stroke. In the current hardware the folded fan is still \SI{30}{\milli\metre} wide
(\SI{12.5}{\square\centi\metre}, 68\,\% of the open area), which makes the spatial asymmetry weak. A folded
width of \SI{12}{\milli\metre} (\SI{5.9}{\square\centi\metre}) is assumed for the optimized stroke in
\cref{sec:res-drag}. The fan is always oriented normal to the stroke; its opening and closing are treated as
instantaneous switches between two rigid shapes (\cref{sec:cfd-composite}).

\paragraph{Lift-type: fluke.} The fluke (design ``A2'') is a small swept foil with a span of
\SI{60}{\milli\metre}, a root chord of \SI{34}{\milli\metre}, a planform area of \SI{11.2}{\square\centi\metre}
and an aspect ratio of about 3.2. Its section is a thick symmetric profile similar to NACA\,0021. It is mounted
on a pin near its leading edge at the lower end of a straight foot rod, so that it can pitch about a spanwise
axis. On the hardware the pitch is passive, set by the hydrodynamic moment against a leaf spring and two end
stops. In the simulations of this thesis the pitch angle $\theta(t)$ is prescribed; how closely a passive
fluke follows a prescribed schedule is not studied here.

\begin{figure}[t]
  \centering
  \includegraphics[width=\textwidth]{fig_propulsors}
  \caption[Propulsor geometries]{Propulsor geometries used in the CFD, drawn to scale. (a) Open fan paddle.
    (b) Folded fan of the current hardware, \SI{30}{\milli\metre} wide. (c) Folded fan of
    \SI{12}{\milli\metre} width assumed for the optimized stroke. (d) Fluke A2 in top view, with its pin
    axis.}
  \label{fig:propulsors}
\end{figure}

\begin{table}[t]
  \centering
  \caption[Propulsor data]{Propulsor data. Areas are planform areas of the propulsor without the linkage.}
  \label{tab:propulsors}
  \small
  \begin{tabularx}{\textwidth}{@{}l>{\raggedright\arraybackslash}X>{\raggedright\arraybackslash}X@{}}
    \toprule
    & Fan paddle (drag type) & Fluke A2 (lift type) \\
    \midrule
    Shape & \SI{120}{\degree} fan, $r=\SI{42}{\milli\metre}$, on \SI{50}{\milli\metre} rod & swept foil, NACA\,0021-like \\
    Area & \SI{18.5}{\square\centi\metre} open; \SI{12.5}{} or \SI{5.9}{\square\centi\metre} folded (30 or 12\,mm wide) & \SI{11.2}{\square\centi\metre} \\
    Size & fan radius \SI{42}{\milli\metre} & chord \SI{34}{\milli\metre}, span \SI{60}{\milli\metre} \\
    Principle & drag, folding recovery & lift, pitching about pin \\
    Degrees of freedom & $\psi,\phi$ + open/closed & $\psi,\phi$ + pitch $\theta$ \\
    Frequency & \SI{0.8}{\hertz} (robot gait), \SI{1.39}{\hertz} (V3trap) & \SI{2}{\hertz} \\
    \bottomrule
  \end{tabularx}
\end{table}

\section{Leg Motions}
\label{sec:plat-motions}

\Cref{fig:kinematics} shows the motions studied.

\paragraph{Robot gait.} The robot currently swims with a gait taken from video of the robot: period
\SI{1.25}{\second}, $\mathrm{DUTY}=0.5$, a sinusoidal hip sweep of \SI{60}{\degree} and the fan folded to
\SI{30}{\milli\metre} during the recovery. It was designed for walking and serves as the baseline.

\paragraph{Optimized drag-type stroke (V3trap).} The stroke derived in \cref{sec:res-drag} has a period of
\SI{0.72}{\second} and $\mathrm{DUTY}=0.382$. The hip sweeps \SI{50}{\degree} centred on the vertical
(\SIrange{120}{70}{\degree}), so that the paddle moves nearly straight backwards and stays normal to its path.
The hip velocity follows a trapezoidal law with ramps of about \SI{45}{\milli\second}. During the recovery the
parallelogram curls the paddle upwards and the fan is folded to \SI{12}{\milli\metre}. The fan is closed at
$\tau_c=0.33$, the start of the power-stroke deceleration, and opened at $\tau_o=0$, the reversal point at the
end of the recovery (\cref{fig:kinematics}a).

\paragraph{Fluke motion.} The fluke leg keeps the hip fixed and swings the foot rod about the knee, so that
the fluke's pin moves up and down on a flat arc: $\pm\SI{24.9}{\degree}$ at \SI{2}{\hertz}, a vertical stroke
of \SI{39.2}{\milli\metre} and a peak heave speed of about \SI{0.28}{\metre\per\second}
(\cref{fig:kinematics}b). The reference motion, called L0, is taken from the whole-robot Flare Live model
(\cref{sec:cfd-flare}): the joint angles and the fluke pitch of one leg were recorded while the fluke pitch
was controlled to hold a target angle of attack at \SI{0.223}{\metre\per\second}, and phase-averaged over the
recorded cycles. The resulting 75th percentile of $|\alpha|$ at this speed is \SI{19.9}{\degree}. Three
families are derived from it:
\begin{itemize}
  \item \emph{L2, fixed pitch schedule:} the L0 motion unchanged at
        \SIlist{0;0.08;0.15;0.30}{\metre\per\second}, as a propeller is tested at fixed rotational speed and
        varying advance speed. Since $\gamma$ changes with $\U$ (\cref{eq:gammamax}), the angle of attack is
        large at low speed and small at high speed.
  \item \emph{L1, angle-of-attack law:} a pitch schedule computed from the heave speed to hold a target
        $\alpha$ of \SI{10}{\degree} or \SI{28}{\degree} at \SI{0.223}{\metre\per\second}, and \SI{20}{\degree} at
        \SI{0.40}{\metre\per\second}. These cases use slightly different hinge paths from L0 and are therefore
        not a clean angle-of-attack sweep.
  \item \emph{L3, speed-scheduled pitch:} the L0 hinge path with the pitch
        schedule~\eqref{eq:l3}, $\theta_0=\gamma_\text{max}/2$, at \SIlist{0.30;0.40}{\metre\per\second}
        ($\theta_0=\SI{21.6}{\degree}$ and \SI{17.6}{\degree}).
\end{itemize}

\begin{figure}[t]
  \centering
  \includegraphics[width=\textwidth]{fig_kinematics}
  \caption[Leg motions]{Leg motions in the hull frame. (a) Optimized drag-type stroke V3trap: paddle positions
    at 25 instants per cycle, open (thick) and closed (thin), and the path of the linkage end point $E$.
    (b) Fluke motion L0: chord positions over one cycle, downstroke and upstroke, with the hinge path. (c) Joint
    histories: hip angle of V3trap and of the robot gait (about their means), fluke pitch angle and hinge
    heave of L0 (dotted, right axis).}
  \label{fig:kinematics}
\end{figure}

Two differences between the propulsors are inherent in this set-up and are kept in mind throughout. First,
the propulsors differ in area (\SI{18.5}{} versus \SI{11.2}{\square\centi\metre}) and in frequency, so thrust
per leg, thrust per area and thrust per power are all reported. Second, the drag-type stroke has been
optimized in several rounds, whereas L0 is a motion taken from another model and was not optimized for
thrust; L3 is the only systematic improvement of the fluke motion. At rest both propulsors draw a similar
hydrodynamic input power of \SIrange{18}{22}{\milli\watt} per leg, so that the comparison is made at a similar
power budget.

\section{Hardware Prototype}
\label{sec:plat-hw}

The lift-type configuration exists as hardware (\cref{fig:hardware}). All four legs carry printed flukes on a passive pin with
a leaf spring. The firmware drives each leg through seven control points per cycle: the thigh is held fixed, the foot bar swings
by $\pm\SI{25}{\degree}$ about the horizontal, and the legs are coordinated in a trot. At \SI{2}{\hertz} this is close to the L0
motion: the pin travels \SI{38.0}{\milli\metre} vertically (L0: \SI{39.2}{\milli\metre}) at a mean heave speed of about
\SI{0.15}{\metre\per\second} (L0: \SI{0.157}{\metre\per\second}), but its peak heave speed is about 15\,\% lower, because the
control points are joined by straight segments. Unlike in the simulations, the hull is mounted reversed, so the robot swims with its
blunt end first. In the water, BODY2 with flukes swam visibly faster than with its fan paddles and its current gait. A first timed
test is compared with the self-propulsion estimate in \cref{sec:res-hw}. The optimized paddle was not tested, because it requires a
controllable fold and a trapezoidal velocity law that the hardware does not have.

\begin{figure}[t]
  \centering
  \includegraphics[width=0.78\textwidth]{fig_hardware.jpg}
  \caption[BODY2 hardware with flukes]{Hardware prototype of BODY2 with printed fluke feet mounted on the leg linkages; two of the
    fluke feet are shown resting on the table. Each leg is driven by two PTK\,7465W servos.}
  \label{fig:hardware}
\end{figure}
